\documentclass[11pt]{article}
\usepackage[margin=25mm]{geometry}
\usepackage{fontspec,unicode-math}
\setmainfont{Latin Modern Roman}
\setmathfont{Latin Modern Math}
\usepackage{amsmath,amssymb,mathtools}
\usepackage{xurl,hyperref}
\hypersetup{colorlinks=true,urlcolor=blue,linkcolor=blue}
\setcounter{secnumdepth}{0}
\setlength{\parindent}{0pt}
\setlength{\parskip}{5pt}
\setlength{\emergencystretch}{3em}
\newcommand{\sref}[2]{\hyperlink{source:#1:#2}{[S#2]}}
\newcommand{\cataloguescope}[1]{\paragraph{Source question.}#1}
\newcommand{\floor}[1]{\left\lfloor#1\right\rfloor}
\begin{document}
\section{Moment-SOS Convergence for PDE-Constrained Variational Problems}\label{statement}

\subsection{Definitions and mathematical statement}

\paragraph{Shared notation.}Standard number systems are denoted by $\mathbb N,\mathbb Z,\mathbb Q,\mathbb R,\mathbb C$; index ranges, computational models and specialized notation follow the cited problem source and the definitions below.


A PDE-constrained variational problem optimizes a functional over functions satisfying a partial differential equation. Finite elements approximate those functions in finite-dimensional spaces; moment-SOS relaxations then replace finite polynomial constraints by increasingly strong positivity conditions.

A polynomial is a sum of squares (SOS) if it is $\sum_j q_j^2$ for polynomials $q_j$. Moment-SOS hierarchies impose positivity through increasingly large truncated moment matrices and corresponding degree-bounded SOS certificates.
\paragraph{Statement recorded in the supplied source.}
Extend the convergence analysis of finite-element moment-sum-of-squares relaxations beyond variational problems with a unique global minimizer and boundary constraints, to variational problems subject to partial differential equation constraints. \sref{statement}{2}
\paragraph{Scope qualification.}
The PDE class, boundary conditions, admissible function space, and relaxation construction are source-specific. The target explicitly extends beyond unique minimizers.
\paragraph{Source status.}
The supplied catalogue labels this record open. Its source-status metadata was checked against the supplied snapshot on 2026-10-11. This import does not independently establish current open status or mathematical correctness. The catalogue literature assessment is dated 2026-08-22.
\subsection{Short English statement}
Prove convergence of the cited finite-element moment method for the broader PDE-constrained setting.
\subsection{Sources}
\begin{itemize}
\item[\hypertarget{source:statement:1}{S1}] Ulam AI and the UnsolvedMath Contributors, UnsolvedMath: A Curated Collection of Open Mathematics Problems, catalogue entry. Catalogue curation and metadata: CC BY 4.0. The supplied source question is retained; definitions and the English summary are editorial additions.
\url{https://huggingface.co/datasets/ulamai/UnsolvedMath}.
\item[\hypertarget{source:statement:2}{S2}] Polynomial Optimization for Nonlinear Dynamics: Theory, Algorithms and Applications (2025), Oberwolfach Reports, DOI 10.4171/owr/2024/35, pp. 1975--2032
\url{https://doi.org/10.4171/owr/2024/35}.
\textit{Underlying Oberwolfach Reports material retains its source-specific licence, including CC BY-SA 4.0 where applicable; see THIRD\_PARTY\_NOTICES.md.}
\end{itemize}
\paragraph{Preparation and review.}
Original question attributed to the cited source and UnsolvedMath contributors. Codex (GPT-6; exact serving revision not exposed) assisted with definitions, English summaries, source inspection and formatting. No new human mathematical authorship or independent proof verification is asserted. Catalogue identity checked against the complete supplied registry and existing repository statements; known flagged repetitions and obvious identical conjectures excluded. Full mathematical-equivalence/literature review has not been performed. Statement preparation and final review followed contribution prompts C and D; automated checks verify structure and identity, not open status.
\label{end:statement}
\end{document}
